\chapter{Methodology}
\label{chap:method}

% See context/outline.md Ch 3 for paragraph-level plan.


\section{Defining the Task}
\label{sec:defining-task}

Ressursplanlegger began as a bachelor task offered by Admmit, a Norwegian software firm building products for the transport sector. The task description, distributed through NTNU as project OPG29, set out a brief: investigate how Norwegian traffic coordinators plan daily transport assignments, and build a working prototype of a planning platform that could assist them. The team accepted the task and committed to delivering both the artefact and the thesis that surrounds it. From that decision onwards, two roles ran in parallel for two students: building Ressursplanlegger as a functioning Next.js platform, and conducting the research that would defend its design choices.

The brief defined the topic but not the empirical foundation, and that part of the work was arranged by the team rather than by Admmit. Seven traffic coordinators and transport managers were contacted directly by phone over a single day in March 2026. All seven worked at companies that were already Admmit customers, but Admmit did not broker the introductions. The team found the coordinators through Admmit's customer roster, called them on its own initiative, and conducted brief semi-structured interviews lasting roughly fifteen minutes each. The interviewed companies range from a small operator with eight to twenty drivers to a multi-department fleet of around forty-five vehicles, and they use a mix of Timpex, Opter, internal tools, and fully manual planning.

Before any coordinator was on the phone, however, two design choices had already been locked. Admmit specified Human-in-the-Loop (HITL) operation as a non-negotiable property of the artefact from the project's first week: the algorithm proposes a plan; the coordinator inspects, modifies, accepts, or rejects each assignment. Admmit also specified a multi-tenant architecture in which a single deployed instance serves multiple customer companies through configuration. Both decisions originate in Admmit's mandate, not in the interviews. The interviews validated HITL by surfacing the coordinator's insistence on retaining authority over assignments, and they validated the cross-company variety that motivates configurable behaviour, but they did not introduce either property. The origin map presented later in this chapter (\Cref{sec:dsr-dsrm}) labels every major design decision by origin and by interview-validation status, so the reader can separate Admmit-mandated choices from interview-driven ones from designer and domain-knowledge choices.

This origin frames the rest of the methodology. Every design choice traced through the chapters that follow comes from one of three sources: the Admmit mandate that fixed HITL and multi-tenancy at the outset, the dialogue with the seven coordinators that informed the planning surface and the assignment-criteria taxonomy, or designer and domain-knowledge decisions that filled in what neither the mandate nor the interviews specified. HITL operationalises the \textbf{Tillit/kontroll} anchor as four concrete coordinator actions on every algorithm-generated assignment, and the multi-tenant architecture enables \textbf{Tilpasningsdyktighet} by isolating per-customer configuration in a single deployed instance. The remaining sections of this chapter describe the research paradigm under which the artefact was built (\Cref{sec:dsr-dsrm}) and how the interview and analysis work was conducted (\Cref{sec:data-collection,sec:data-analysis}). The eight named iterations that produced the system are documented in \Cref{sec:iterations}, and the framework under which it was evaluated, together with the validity bounds that follow, is set out in \Cref{sec:evaluation-framework,sec:validity}.


\section{Method Theory --- DSR + DSRM Applied}
\label{sec:dsr-dsrm}

\Cref{sec:defining-task} mapped where each design choice originated; this section names the research paradigm under which those choices were operationalised. The DSR foundation and the positivist and interpretive alternatives that frame it are set out in \Cref{sec:dsr}, and that comparison is not repeated here. A purely positivist design would require deployment-scale measurement of overtime and idle-time effects in Norwegian transport companies, which the present project's stage does not support, while a purely interpretive design would yield rich description of coordinator practice without producing the runnable platform the bachelor task required. DSR is the synthesis the project requires: \textcite{hevner2004design} characterise it as a paradigm in which knowledge of a problem and its solution is acquired through the building and application of a purposeful artefact \parencite[p.~82]{hevner2004design}, and Ressursplanlegger is precisely such an artefact.

The structured operationalisation used in this thesis is \textcite{peffers2007dsrm}'s Design Science Research Methodology (DSRM), which renders DSR as six nominally sequential activities: problem identification and motivation; objectives for a solution; design and development; demonstration; evaluation; and communication \parencite[p.~46]{peffers2007dsrm}. The activities are nominally sequential rather than strictly so; \textcite{peffers2007dsrm} note explicitly that researchers are not expected to proceed in order from activity~1 through activity~6 \parencite[p.~56]{peffers2007dsrm}. The six activities below describe what was done in this project at each step.

\begin{description}
  \item[Problem identification and motivation.] The seven coordinator interviews surfaced a resource-utilization visibility gap: overtime is handled reactively, idle time between assignments is invisible, and load balancing depends on coordinator memory rather than any structured overview. Admmit's bachelor task framed the same gap from the owner side, asking for an algorithm-assisted planning platform that could make utilisation legible and adjustable. \Cref{sec:interview-findings} describes the empirical gap in detail; \Cref{sec:effektivitet} discusses it under the Effektivitet anchor.
  \item[Objectives for a solution.] The objectives are operationalised under the three locked anchor concepts. \textbf{Effektivitet} is realised by the artefact reducing overtime, reducing idle time between assignments, and reducing uneven load between drivers. \textbf{Tillit/kontroll} is operationalised as the coordinator's authority to inspect, modify, accept, or reject any algorithm-generated assignment. \textbf{Tilpasningsdyktighet} is delivered by configurable constraint weights and per-company assignment rules within a single deployed instance that can serve companies with different operational rules.
  \item[Design and development.] The artefact is an algorithm-assisted planning platform built on Next.js, tRPC, and PostgreSQL, with a multi-engine solver layer (a greedy heuristic, the OR-Tools CP-SAT constraint solver, and the Timefold metaheuristic) invoked as subprocesses, and a drag-and-drop timeline through which the coordinator inspects, modifies, accepts, or rejects each assignment. The platform was built across eight named iterations, each labelled by origin (Admmit-mandated, interview-driven, or designer / domain-knowledge), described in \Cref{sec:iterations}.
  \item[Demonstration.] Demonstration takes the form of a multi-engine benchmark on synthetic datasets sized to span realistic constraint combinations across the seven interviewed companies (small fleets of around ten drivers through to multi-department operations of around forty-five vehicles), used to show how the three solver families behave under the same problem instances. The dataset design and benchmark protocol are set out in \Cref{sec:evaluation-framework}.
  \item[Evaluation.] What the artefact is \emph{validated against} is \emph{how} solver approaches compare under realistic constraint combinations, not \emph{whether} the visibility gap is real or \emph{whether} the HITL stance is necessary. The benchmark is a methodologically independent validation instrument for solver performance, and the requirements traceability matrix records coverage of functional and non-functional requirements by recording, for each requirement, whether it was implemented and how its behaviour is checked. What is \emph{not validated} in this thesis is forwarded to \Cref{sec:limitations}: the artefact is not deployed in production, real-world utilisation gains are not measured, and the override flow is not user-tested with coordinators. The validation-versus-evaluation distinction that bounds these claims is set out under DSR theory in \Cref{sec:dsr} and applied at the level of the benchmark in \Cref{sec:evaluation-framework}.
  \item[Communication.] The research is communicated through two channels. The first is this thesis, addressed to the academic audience at NTNU and presenting Ressursplanlegger as an instantiation in DSR's sense alongside the design knowledge embedded in its multi-engine and HITL choices. The second is a written recommendation to Admmit covering deployment readiness, including the requirements not yet validated through user testing (\Cref{sec:limitations}), and the per-company configurability surface through which Tilpasningsdyktighet is delivered.
\end{description}

The mapping above places the project on Peffers' nominal sequence, but the actual ordering reflected the project's entry point rather than the sequence. \textcite{peffers2007dsrm} note four entry points into the DSRM cycle, and explicitly identify objective-centered initiation, in which an industry or research need fixes key objectives before activity~1 is complete, as a legitimate starting point \parencite[p.~56]{peffers2007dsrm}. Ressursplanlegger fits that pattern: Admmit's mandate fixed Tillit/kontroll and Tilpasningsdyktighet before activity~1 was complete, and the eight iterations of activity~3 ran in parallel with later refinements of the objectives.


\section{Data Collection}
\label{sec:data-collection}

The Problem-Identification activity laid out in \Cref{sec:dsr-dsrm} drew its empirical input from a single instrument: brief semi-structured interviews with traffic coordinators and transport managers across seven Norwegian transport companies. The semi-structured form was chosen because the activity needed data open enough to surface unanticipated themes yet structured enough that the same five guide topics could be compared across the seven calls. The visibility gap and the operator-vs-owner asymmetry both surfaced this way rather than from any prepared question. A fully structured questionnaire would have closed off unanticipated themes; a fully open conversation would have lost cross-case comparability. The methodological framing for the analytic step that turned these interviews into themes is set out in \Cref{sec:data-analysis}.

With the form of the instrument fixed, the next decision was who to consult. The seven coordinators were selected purposively from Admmit's customer roster rather than recruited through a public call, because the project required coordinators in the planning role at companies where algorithm-assisted planning is plausibly relevant, and Admmit's customer base supplied that population directly. The companies span small operators with around eight to twenty drivers through to a multi-department fleet of around forty-five vehicles, range geographically from Bergen and Mo i Rana to Lakselv and Tana, and cover the live system landscape: Timpex, Opter, internal or custom tools, and fully manual planning. The self-selection that follows from the Admmit-customer recruitment route is a known bound on the empirical foundation and is named as L2 in \Cref{sec:limitations}.

Within those seven calls, the same five guide topics anchored every conversation: current tools and workflow, assignment criteria, sick-leave handling, attitudes to automation, and adoption criteria. Each topic opened with a broad question and was followed up with prompts only where the coordinator's first answer left an obvious gap; this kept the calls comparable on topic coverage while preserving the open-question structure the form requires. The full topic guide is documented in the project's working notes.

The seven calls were conducted by phone over a single day in March 2026, in Norwegian, with each call lasting approximately fifteen minutes for a cumulative interview time of roughly 105 minutes. The single-day window kept the interviewer's framing consistent across calls and avoided the drift that comes from conducting interviews weeks apart, but it also constrained how deeply any one topic could be probed within a single call; the brief duration is named alongside the sample size as L1 in \Cref{sec:limitations}.

Recordings were transcribed automatically and the segments used for analysis were manually corrected against the audio. Automated transcription was chosen because the cumulative audio of around 105 minutes made fully manual transcription disproportionate to the analytic depth the consultations could support.

Beyond the instrument and its mechanics, the consultation framing has ethical consequences worth stating directly. The project was conducted as an industry consultation under Admmit's bachelor task brief and NTNU institutional supervision rather than as formal research, and on that basis no Sikt (Norway's national agency for shared services in research and education) registration was sought. The conversations covered work practices, system use, and assignment criteria only, and no special-category personal data was collected. Consent was implicit through participants agreeing to take the call and answer questions; no formal consent form was administered, and recording proceeded without an explicit pre-recording consent step. No formal anonymisation protocol was applied. Only Timpex and Opter, the two commercial vendors named in the consultation pool, are named in the thesis, while the seven interviewed companies are referenced generically by size, sector, and system maturity. Recordings and transcriptions are stored locally on the project author's device, accessible only to the author and supervisor, and will be deleted after submission and grading. The researcher--developer overlap and the informal study procedures that follow from this consultation framing are named together as L3 in \Cref{sec:limitations}.


\section{Data Analysis}
\label{sec:data-analysis}

The seven interviews described in \Cref{sec:data-collection} produced roughly 105 minutes of audio across the five guide topics that anchored each call (current tools and workflow, assignment criteria, sick-leave handling, attitudes to automation, and adoption criteria). The analytic instrument applied to that material was thematic analysis as set out by \textcite{braun2006thematic}, chosen for its theoretical flexibility within the DSR frame \parencite[p.~81]{braun2006thematic} and its accessibility to a small research team conducting analysis as one component of a larger DSR Problem-Identification activity \parencite[p.~77]{braun2006thematic}, conducted across their six phases of familiarisation, initial coding, theme search, theme review, theme definition, and report production \parencite[p.~87]{braun2006thematic}. Three positioning choices situate the analysis within the room \textcite{braun2006thematic} explicitly require analysts to occupy \parencite[pp.~81--82]{braun2006thematic}: it was \emph{inductive}, in that no codebook was drawn from a prior hypothesis and the visibility gap and operator-vs-owner asymmetry were named only after coding; \emph{semantic}, in that codes worked at the level of coordinators' explicit statements about tools, criteria, sick-leave, automation, and adoption rather than at a latent discursive level; and \emph{essentialist/realist}, in that coordinator accounts were treated as reports of real working practice usable as input to DSR Problem Identification (\Cref{sec:dsr-dsrm}). The analysis identified themes across the seven transcripts that elaborated and cut across the five guide topics (among them the resource-utilization visibility gap, the operator-vs-owner asymmetry, and automation scepticism with an override preference); following \textcite[Table~2, criterion~15]{braun2006thematic}, those themes were named through the coding process rather than treated as having surfaced unaided. Braun and Clarke's own examples are drawn from psychology and gender studies; the application here is to brief semi-structured consultations with traffic coordinators inside a DSR Problem-Identification activity, not to standalone qualitative research, and the six-phase guide travels across that domain gap by treating each transcript as a data item and the seven transcripts together as the data set.

The output of the analysis fed three downstream sections. The named themes are reported in \Cref{sec:interview-findings}. Each theme was then translated into one or more functional requirements prioritised under MoSCoW (Must, Should, Could, Won't), chosen for its clear stakeholder-negotiable boundary between in-scope and out-of-scope requirements, recorded in the functional-requirements register (IDs FK-01 through FK-42) and reported in \Cref{sec:requirements}: the visibility-gap theme produced the timeline view and conflict-detection requirements at Must, the assignment-criteria theme produced the hard- and soft-constraint requirements at Must and Should respectively, and the automation-attitudes theme produced the override and weight-configuration requirements at Must and Should. A fit/gap analysis --- a side-by-side comparison of what existing systems cover and what they do not --- was then run against Timpex, Opter, and the internal or manual tools named in the consultation pool, producing five fit items where existing systems already meet the need and fourteen gap items where they do not, reported in \Cref{sec:fit-gap}.

The analysis carries the limitations of its inputs and its frame, which are catalogued in \Cref{sec:limitations} rather than re-explained here: the single-day collection window is named as L1 alongside sample size, and the recruitment of all seven coordinators from Admmit's customer roster is named as L2 as self-selection bias. The themes are accordingly bounded to the seven interviewed companies and are not a basis for generalising to Norwegian transport more broadly.


\section{Iterative Development Process}
\label{sec:iterations}

The artefact described in \Cref{sec:dsr-dsrm} was built across eight named iterations, told here as a connected narrative rather than as a sprint diary. Each iteration follows the same pattern: what was tried, why it was tried, what happened (positive results and limitations), what was learned, and what motivated the next step. The opening iteration establishes the data substrate (\Cref{sec:iter-schema}); iterations two through four work outward from a single greedy heuristic to a multi-engine solver layer (\Cref{sec:iter-greedy,sec:iter-cpsat,sec:iter-multiengine}); iterations five through seven build the operator surface that lets the coordinator inspect, modify, accept, or reject every algorithm-generated assignment (\Cref{sec:iter-hitl,sec:iter-deviations,sec:iter-tradeoff}); and the eighth iteration exposes per-company configurability of the soft-constraint weight space (\Cref{sec:iter-weights}). Each iteration carries an inline origin label per the origin map introduced in \Cref{sec:defining-task}: Admmit-mandate, interview-validated, designer-technical, or a mixture. Iteration boundaries are reconstructed retrospectively from project artefacts and decision points; precise dates and ordering are approximate. The narrative below describes how the artefact was built; the framework under which it was tested follows in \Cref{sec:evaluation-framework}.

\subsection{Schema and Timeline Scaffolding}
\label{sec:iter-schema}

\begin{description}
  \item[Origin.] Admmit-mandate (multi-tenant architecture) combined with designer-technical exploration of the data model and a read-only timeline UI.
  \item[Tried.] A core entity model covering employees, vehicles, assignments, certifications, weekly work schedules, and time-off; multi-tenant data isolation enforced at the data layer with every query scoped by \texttt{companyId}; and a read-only Gantt-style timeline that rendered existing assignments without any optimisation step.
  \item[Why.] Later solver iterations need a structured substrate to operate over, and the multi-tenant boundary must be enforced at the data layer because Admmit's customer structure means each transport company would consume the artefact independently; bolting tenancy on later would risk cross-tenant leakage on every new query path.
  \item[What happened.] The substrate was adequate for the simplest assignment heuristic but required revision once constraint-solver formalisation surfaced gaps in the data model. Explicit time-off windows, certificate-expiry as a queryable date, and per-day work-schedule entries became necessary only after the heuristic exposed what the solver would need to reason about.
  \item[Learned.] Schema-first discipline reduces rework but cannot fully anticipate solver requirements. The substrate is itself an iteration, not a one-shot setup.
  \item[Next.] With the substrate stable enough to support a real assignment step, the next iteration attempted the simplest possible assignment heuristic.
\end{description}

\subsection{Single-engine Baseline (greedy)}
\label{sec:iter-greedy}

\begin{description}
  \item[Origin.] Designer-technical exploration.
  \item[Tried.] A priority-sorted greedy first-fit assignment heuristic. Assignments were sorted by priority, walked through in order, and each was given the first feasible driver and vehicle combination found.
  \item[Why.] An instant baseline that produces a feasible plan provides a reference point for measuring richer-solver gains and gives the coordinator immediate UI feedback while drag-and-drop edits happen.
  \item[What happened.] The heuristic worked on small instances of fewer than thirty assignments. On larger instances, priority-sorted greedy committed early to locally optimal but globally suboptimal choices, blocking later high-priority assignments and producing avoidable idle time and overtime.
  \item[Learned.] The greedy ceiling is structural rather than a tuning issue; reordering the priority comparator cannot break out of it because the heuristic lacks both constraint propagation and a global view of the day.
  \item[Next.] Introduce a complete solver that can reason about all assignments simultaneously.
\end{description}

\subsection{Constraint Generalisation (CP-SAT)}
\label{sec:iter-cpsat}

\begin{description}
  \item[Origin.] Designer-technical exploration.
  \item[Tried.] The Google OR-Tools CP-SAT (Constraint Programming Satisfiability) solver with explicit hard constraints (competencies, availability, no double-booking, vehicle type) and weighted soft constraints (workload balance, vehicle preference, transitions, priority, employee preferences) combined under an objective function maximising scheduled assignments minus weighted soft-constraint violations.
  \item[Why.] Constraint programming is a known fit for assignment problems with rich interaction between hard and soft constraints, and CP-SAT specifically returns near-optimal solutions within a configurable time budget, a property that matters for an HITL surface where the coordinator is waiting on the solver.
  \item[What happened.] The solver returned near-optimal solutions on small to medium instances within a thirty-to-sixty-second time limit. Quality degraded on larger instances when the time limit bound: the solver returned the best feasible solution found rather than the optimum. Modelling complexity was higher than for the heuristic, and several iterations were needed to get the hard-versus-soft mapping correct.
  \item[Learned.] CP-SAT covers the medium-instance gap that the greedy heuristic could not, but introduces a scaling boundary that motivates a third engine for larger instances.
  \item[Next.] Add a third engine for large instances and a mechanism for comparing the three under matched conditions.
\end{description}

\subsection{Multi-engine Selection and Benchmarking}
\label{sec:iter-multiengine}

\begin{description}
  \item[Origin.] Designer-technical exploration. Operationalises the methodologically independent test framing carried through to \Cref{sec:evaluation-framework}.
  \item[Tried.] A pluggable solver registry with three engines: the greedy heuristic as instant baseline, OR-Tools CP-SAT as a complete-with-time-budget solver, and the Timefold metaheuristic (tabu search, simulated annealing, late-acceptance hill climbing) for large-instance scaling. A benchmarking framework with synthetic small, medium, and large datasets ran the same problem instances against all three engines, with engine-version snapshots recorded for reproducibility.
  \item[Why.] Under realistic constraint combinations, the open question is which solver approach best meets Effektivitet on the metrics that the framework in \Cref{sec:evaluation-framework} operationalises (scheduled-assignment count, soft-constraint penalty, runtime within budget). That is a how-question, independent of whether the visibility gap is real (the interviews surfaced this) or whether HITL is necessary (Admmit fixed this from the project's first week). The multi-engine architecture is the instrument designed to answer that how-question.
  \item[What happened.] Subprocess plumbing was complex: Timefold runs as a separate JVM process, CP-SAT inside a Python process, and the greedy solver inside the TypeScript runtime. Early benchmark results were inconsistent across runs until engine-version snapshots were added; the inconsistencies traced to silent solver-library updates between runs rather than to the problem instances themselves.
  \item[Learned.] The multi-engine architecture is a comparative test of \emph{how}, not \emph{whether}. It measures Effektivitet across solver families rather than validating that the underlying problem is worth solving; that framing returns in \Cref{sec:evaluation-framework} and at the discussion level under the Effektivitet anchor.
  \item[Next.] With the computation infrastructure stable, build the human-in-the-loop surface that would let the coordinator act on the plans the solvers produced.
\end{description}

\subsection{HITL Surface}
\label{sec:iter-hitl}

\begin{description}
  \item[Origin.] Mixed: Admmit-mandate (HITL as a non-negotiable principle from the project's first week) combined with interview-validated direction (the drag-and-drop interaction style was confirmed by coordinator consensus that the system should suggest a plan the coordinator can correct, not replace the coordinator).
  \item[Tried.] A Gantt-style drag-and-drop assignment UI, an explicit inspect / modify / accept / reject control on every algorithm-generated suggestion, and a thin post-hoc constraint check that surfaced basic feasibility violations such as over-booking and missing competence.
  \item[Why.] HITL was Admmit's requirement from project start, and the interviews validated the necessity through a coordinator consensus that automated planning must be correctable rather than authoritative. Without an interactive override path, the coordinator would have authority in principle but no way to exercise it on the timeline they actually work with.
  \item[What happened.] Drag-and-drop worked. State consistency between optimistic UI updates and server mutations occasionally produced stale conflict displays, where the timeline showed a violation the server had already resolved, and required a refresh hook on assignment updates. The thin post-hoc check caught the most common conflicts but was clearly insufficient for the range of edge cases the coordinator would encounter.
  \item[Learned.] The HITL surface needs richer real-time feedback than a thin post-hoc check provides. Coordinator authority over plan content is meaningful only when the system surfaces what to act on; an empty timeline the coordinator edits blindly delivers Tillit/kontroll only formally.
  \item[Next.] Expand the post-hoc check into a full deviation taxonomy that materialises the constraints the coordinator would otherwise have to remember.
\end{description}

\subsection{Real-time Deviation Detection}
\label{sec:iter-deviations}

\begin{description}
  \item[Origin.] Designer and domain knowledge. The deviation categories extend pain points named in interviews around tacit knowledge and overtime, but the specific category set was not asked about in the interview guide.
  \item[Tried.] A post-generation conflict scanner across multiple deviation categories: overbooking, overtime violations, missing competencies, expired certifications, rest-period violations, vehicle maintenance conflicts, night-work conflicts, and assignments to unavailable employees. Each deviation carries a severity level and a status workflow (open, acknowledged, resolved). Deviations appear inline on the planning timeline and on a dedicated deviations view.
  \item[Why.] A meaningful HITL surface requires the system to surface what the coordinator should override toward, not just an empty timeline they edit blindly. The deviation taxonomy is the mechanism through which constraints kept in the coordinator's head become structured conflict data on the timeline.
  \item[What happened.] False positives appeared at edge-case interactions between weekly work schedules and time-off windows: week boundaries and mid-shift absences produced noisy alerts. Rest-period violations were inherently noisy because the regulation admits multiple definitions (the forty-eight-hour weekly average, the sixty-hour weekly maximum, the daily thirteen-hour limit), and a deviation flagged under one definition could be compliant under another.
  \item[Learned.] Deviation detection materialises tacit knowledge as structured conflict data, but tuning thresholds is itself a coordination problem the system cannot fully resolve. The coordinator decides when a flagged deviation is real, which keeps the override authority where the HITL stance places it.
  \item[Next.] With the planning surface and deviation feedback in place, expose the user-facing planning-time control through which the coordinator chooses what kind of plan to ask for.
\end{description}

\subsection{User-controlled Plan-time vs Plan-quality Tradeoff}
\label{sec:iter-tradeoff}

\begin{description}
  \item[Origin.] Designer-technical, framed as a Tillit/kontroll refinement that gives the coordinator authority over the planning process itself, not only over plan outcomes.
  \item[Tried.] A coordinator-facing solver-time-budget control. The coordinator selects either \emph{fast / less accurate} (the greedy heuristic with a short time limit) or \emph{slower / more accurate} (the CP-SAT solver with a longer time limit). The time-quality tradeoff is exposed as a deliberate operator choice rather than an internal default.
  \item[Why.] Different planning situations carry different urgency-quality requirements: a sick-leave replanning event under time pressure differs from next week's planning. Coordinator authority over how the plan is generated is a Tillit/kontroll dimension distinct from authority over plan content; the override flow handles the latter, the time-budget control adds the former.
  \item[What happened.] The choice surface itself raised the question of what the user is choosing between. Bare time labels (``ten seconds'' versus ``sixty seconds'') were not legible to a coordinator who does not think in solver-runtime budgets; descriptive labels (``fast suggestion'' versus ``best-effort plan'') were clearer but smuggled in claims about plan quality the system cannot verify per instance. The compromise settled on descriptive labels with a tooltip explaining what the toggle changes.
  \item[Learned.] Exposing a tradeoff in the interface requires the interface itself to explain it, which connects this control to the explanation-as-interface treatment of HITL set out in \Cref{sec:hitl}. The time-budget toggle is a small case of the same principle: the coordinator can act on the tradeoff only when the interface makes the tradeoff legible.
  \item[Next.] Move from operator authority over the planning process to per-company configurability of the planning rules themselves: the soft-constraint weight space.
\end{description}

\subsection{Configurable Soft-constraint Weights}
\label{sec:iter-weights}

\begin{description}
  \item[Origin.] Designer and architectural. Operationalises the Tilpasningsdyktighet anchor that the project committed to from the start through the Admmit-mandated multi-tenant architecture. The specific weight set was chosen by the designers, not interview-validated as the right set.
  \item[Tried.] A per-company weight configuration UI for soft constraints: workload balance, vehicle preference, transitions between consecutive assignments, priority, and employee preferences. Each weight flows into the CP-SAT objective function so that the same engine produces different plans for different companies without code changes.
  \item[Why.] The interviews confirmed that assignment criteria differ materially across companies, with some weighting workload balance highly and others prioritising vehicle--driver continuity or route familiarity. Tilpasningsdyktighet means the same artefact must serve materially different operational rules within a single deployed instance.
  \item[What happened.] User-chosen weight combinations were hard to validate against degenerate solutions. Weights summing to zero produced unstable rankings, and one weight dominating to the exclusion of others produced plans that satisfied a single criterion at the cost of all others. The UX challenge was explaining what each weight actually means to a coordinator who does not work in optimisation vocabulary.
  \item[Learned.] Configurability for soft weights is a partial Tilpasningsdyktighet realisation. Hard constraints, status workflow definitions, and per-company assignment-rule taxonomies remain hard-coded; broader configurability of those layers is acknowledged as L4-adjacent in \Cref{sec:limitations} and forwarded to future work in \Cref{sec:future-work}.
  \item[Next.] With eight iterations of artefact in place, the methodological question becomes how to test whether the artefact meets the locked anchors. That question is answered in \Cref{sec:evaluation-framework}.
\end{description}


\section{Evaluation Framework}
\label{sec:evaluation-framework}

The eight iterations in \Cref{sec:iterations} produced an artefact whose claims now need a framework to test them against. That framework is asymmetric on purpose, following the position carried through \Cref{sec:iter-multiengine}: Ressursplanlegger embodies two claims it does not test, and tests one claim it does not embody. The artefact embodies the position that the visibility gap is real, which the seven coordinator interviews surfaced (\Cref{sec:data-collection}). It also embodies the position that human-in-the-loop operation is the right operator stance, which Admmit fixed in the project's first week (\Cref{sec:defining-task}). What the artefact tests, through the multi-engine architecture and the benchmark protocol described below, is \emph{how} solver approaches compare under realistic constraint combinations, not \emph{whether} the underlying problem is worth solving. The validate-versus-evaluate distinction that bounds these claims is set out under DSR theory in \Cref{sec:dsr} and is not repeated here; the benchmark is a validation instrument in Wieringa's sense (see \Cref{sec:dsr}), applied at the level of solver behaviour.

This how-not-whether stance carries through to how each anchor is exercised, and the mapping is asymmetric. \textbf{Effektivitet} has the tightest binding: the multi-engine benchmark is the validation instrument designed for it, comparing scheduled-assignment count, soft-constraint penalty, and runtime-within-budget across the greedy heuristic, OR-Tools CP-SAT, and Timefold metaheuristic on identical problem instances. \textbf{Tilpasningsdyktighet} is exercised only indirectly: the requirements traceability matrix records that the per-company configurable weight surface (\Cref{sec:iter-weights}) is implemented and verified during development, but no benchmark compares the artefact's behaviour across two genuinely different operational rulebooks. \textbf{Tillit/kontroll} sits in the same partial-coverage position as Tilpasningsdyktighet, with one acknowledged gap: the inspect / modify / accept / reject flow is implemented, but it has not been user-tested with traffic coordinators, and that gap is forwarded to L8 in \Cref{sec:limitations}.

The benchmark runs against three synthetic datasets whose sizes were chosen to exercise solver behaviour at three operational scales. The small dataset (five employees, three vehicles, twenty assignments) is a sanity-scale instance on which every engine should finish in well under a second; the medium dataset (twenty employees, ten vehicles, one hundred assignments) exercises the overlap and workload-balance terms without saturating the search; the large dataset (fifty employees, twenty-five vehicles, five hundred assignments) is a stress-scale instance designed to push CP-SAT toward its sixty-second time budget and probe metaheuristic behaviour. The choice of synthetic over production data follows from two concrete constraints rather than from methodological preference: no transport company has operated the artefact, so production data does not exist for this codebase, and the operational data the seven interviewed companies hold falls under the data-protection bounds the consultation framing in \Cref{sec:data-collection} chose not to cross. The price the framework pays for that choice is that benchmark numbers cannot be projected onto real Norwegian-transport operational data without further validation, a bound named as L5 in \Cref{sec:limitations}.

Alongside the benchmark, the requirements traceability matrix functions as the coverage check rather than as a quality validation. For each functional requirement in the FK-01 to FK-42 range, the matrix records two things: whether it is implemented in the deployed codebase, and whether its behaviour was verified during development. ``Verified during development'' is a narrower claim than passing an automated test suite, and it is the honest one to make: thirty-six of the thirty-seven in-scope functional requirements are implemented and developer-verified, one is partial (FK-30, driving and rest-time status), and the five out-of-scope requirements (FK-38 through FK-42) are documented as future work. The absence of a formal automated test suite is itself a non-functional gap and is named as L9 in \Cref{sec:limitations}; the traceability matrix's role is to make coverage legible, not to substitute for the testing that L9 acknowledges is missing.

Three claims about the artefact are deliberately outside what this framework tests, and each forwards to a specific limitation in \Cref{sec:limitations}. The artefact is not deployed in production, so claims about Effektivitet gains under operational time pressure cannot be made from these results: that gap is L6. There is no empirical comparison against Timpex, Opter, or the internal and custom tools the seven interviewed companies use, so claims that Ressursplanlegger is better than alternatives in the same market segment are not supported here: that gap is L7. The inspect / modify / accept / reject flow that operationalises Tillit/kontroll has not been user-tested with traffic coordinators, so claims about its comprehensibility or its behaviour under coordinator workload cannot be made: that gap is L8. The framework that follows in \Cref{sec:validity} addresses the research's epistemic limits rather than the artefact's; this section ends where validation ends, and \Cref{sec:limitations} carries forward what neither addresses.


\section{Validity and Reliability}
\label{sec:validity}
